% !TeX program = xelatex
% 所有数据与给定区间均为虚构；MIT License。
\documentclass[fontset=fandol]{ctexart}
\usepackage[margin=18mm]{geometry}
\usepackage{gradbars}
\usepackage{gradbars-contrib-pairs}
\gradbarsdeclarepaircolumn{service-pair}{frontend=dumbbell,
 options={min=0,max=100,precision=1,unit={\%}}}
\gradbarsdeclarepaircolumn{service-span}{frontend=range,
 options={min=0,max=100,precision=1,unit={\%}}}
\newcommand\servicetable[2]{%
 \begin{center}
 \begin{gradtable}[style=paired-report,theme=#1,width=\linewidth,
   graphic width=22mm,label width=25mm]
  \gradcolumn[width=25mm]{name}{服务项目}
  \gradcolumn[type=service-pair,options={#2}]{pair}{当前 / 参考}
  \gradcolumn[type=service-span]{span}{给定区间}
  \gradrow{预约办理,66/82,60/88}
  \gradrow{上门服务,0/12,0/20}
  \gradrow{咨询服务,84/NA,NA}
  \gradrow{社区活动,NA/76,65/85}
  \gradrow{待复核,NA,NA}
 \end{gradtable}
 \end{center}}
\begin{document}
\section*{社区服务覆盖率的配对对照}
示例数据仅用于演示，不代表真实实验结果。
配对输入顺序为参考值/当前值，原生标签显示顺序为当前值/参考值。
给定区间由使用者提供，不表示推断得到的置信区间。
\servicetable{accessible}{}
\subsection*{黑白交付}
\servicetable{mono}{color=black,compare color=black!65,track color=black!25}
\subsection*{只使用公开接口的记录消费者}
图后数字是当前值减参考值的原值差，再格式化到一位小数；
n/a 表示这组数据不完整。
\par\medskip
\gradpairgraphic[max=100,width=30mm,precision=1]{review-delta}{66}{82}\par\medskip
\gradpairgraphic[max=100,width=30mm,precision=1]{review-delta}{NA}{0}
\end{document}
